\documentclass[11pt]{article}

\usepackage[margin=1in]{geometry}
\usepackage{amsmath,amssymb,amsthm}
\usepackage{booktabs}
\usepackage[hidelinks]{hyperref}
\hypersetup{
  pdftitle={The Width-Five q-Fibonomial Coefficients Are Unimodal},
  pdfauthor={Hainan Zhao}
}

\newtheorem{theorem}{Theorem}
\newtheorem{lemma}[theorem]{Lemma}
\newcommand{\qint}[1]{[#1]_q}
\newcommand{\qfib}[2]{\left[\!\begin{matrix}#1\\#2\end{matrix}\!\right]_{\!\mathcal F}}

\title{The Width-Five \(q\)-Fibonomial Coefficients Are Unimodal}
\author{Hainan Zhao}
\date{August 2026}

\begin{document}
\maketitle

\begin{abstract}
Bergeron, Ceballos, and K\"ustner conjectured that every
\(q\)-Fibonomial coefficient is unimodal.  The cases of width at most four
are known.  We prove the next infinite family: for every integer \(m\geq1\),
the polynomial \(\qfib{m+5}{5}\) is unimodal.  Its first coefficient
difference reduces below the midpoint to six translates of the restricted
partition function for parts \(1,2,3,5\).  That function is a cubic
quasipolynomial whose nonconstant terms are independent of the residue
class.  Uniform error bounds prove positivity for \(m\geq8\), and the seven
remaining cases are checked exactly.
\end{abstract}

\section{Introduction}

Let \(F_0=0,F_1=1\), and \(F_{j+2}=F_{j+1}+F_j\).  For a positive integer
\(r\), write
\[
 [r]_q=1+q+\cdots+q^{r-1},\qquad
 [F_r]^!_q=\prod_{j=1}^r[F_j]_q.
\]
The \(q\)-Fibonomial coefficient is
\[
 \qfib{m+n}{n}=\frac{[F_{m+n}]^!_q}{[F_m]^!_q[F_n]^!_q}.
\]
Bergeron, Ceballos, and K\"ustner proved that this is a polynomial with
nonnegative integer coefficients and conjectured that it is unimodal
\cite{BCK}.  Connelly, Ito, Martinez, Shevchenko, and Yang proved the cases
\(n\leq3\) \cite{CIMSY}; the case \(n=4\) was subsequently proved by the
restricted-partition method used here \cite{ZhaoWidthFour}.

\begin{theorem}\label{thm:main}
For every integer \(m\geq1\), the \(q\)-Fibonomial coefficient
\(\qfib{m+5}{5}\) is unimodal.
\end{theorem}

The argument is algebraic.  The path--domino theorem of \cite{BCK} is used
only for polynomiality and nonnegativity.  We make no claim here about width
six, the full two-parameter conjecture, log-concavity, or a combinatorial
chain decomposition.

There is also an aligned-center criterion for products with several
lacunar brackets: after pairing disjoint factors whenever a spacer step
divides an ordinary bracket length, the remaining spacer increments may be
allocated among the remaining ordinary brackets, provided every residual
length stays positive.  Applied to the three steps (2,3,5) in (2), that
criterion proves exactly 18 of the 60 stable residue classes modulo 60.  It
does not reach 24 further classes that admit such a three-spacer
decomposition, nor the 18 classes without one.  Thus it does not subsume
Theorem~\ref{thm:main}; the restricted-partition argument below treats all
classes uniformly.

\section{The restricted-partition kernel}

For an integer \(t\), define
\[
 p(t)=[q^t]\frac{1}{(1-q)(1-q^2)(1-q^3)(1-q^5)},
 \qquad p(t)=0\quad(t<0),
\]
and put
\[
 Q(t)=\frac{t^3}{180}+\frac{11t^2}{120}+\frac{9t}{20}.
\]

\begin{lemma}\label{lem:kernel}
For \(t\geq0\),
\[
 p(t)=Q(t)+\rho_{t\bmod30},
\]
where the values of \(360\rho_r\), for \(r=0,1,\ldots,29\), are
\[
\begin{array}{c|rrrrrrrrrr}
r&0&1&2&3&4&5&6&7&8&9\\ \hline
360\rho_r&360&163&248&243&136&275&288&163&248&171\\[2pt]
r&10&11&12&13&14&15&16&17&18&19\\ \hline
360\rho_r&280&203&288&163&176&315&208&203&288&91\\[2pt]
r&20&21&22&23&24&25&26&27&28&29\\ \hline
360\rho_r&320&243&208&203&216&235&248&243&208&131
\end{array}
\]
In particular,
\[
 Q(t)+\frac{91}{360}\leq p(t)\leq Q(t)+1. \tag{1}
\]
\end{lemma}

\begin{proof}
The generating function shows that \(p(t)\) counts solutions of
\(x+2y+3z+5w=t\) in nonnegative integers.  Summing first over \(w\) gives
\[
 p(t)=\sum_{w=0}^{\lfloor t/5\rfloor}
 \left\lfloor\frac{(t-5w)^2+6(t-5w)+12}{12}\right\rfloor,
\]
using the period-six formula for partitions with parts \(1,2,3\).
Write \(t=30u+r\), split \(w\) by its residue modulo six, and sum the
resulting quadratic progressions.  The cubic, quadratic, and linear terms
are independent of \(r\) and give \(Q(t)\); the constant terms are precisely
the displayed table.  Its minimum and maximum give (1).
\end{proof}

\section{The midpoint difference}

Fix \(m\geq1\) and set
\[
 a=F_{m+1},\qquad b=F_{m+2}.
\]
The five numerator lengths are
\(a,b,a+b,a+2b,2a+3b\), so
\[
 W_m(q):=\qfib{m+5}{5}
 =\frac{[a]_q[b]_q[a+b]_q[a+2b]_q[2a+3b]_q}
 {[2]_q[3]_q[5]_q}. \tag{2}
\]
Its degree is
\[
 D=5a+7b-12. \tag{3}
\]
As in the width-four case, q-reciprocity, polynomiality, and the constant
term one show that \(W_m\) is symmetric of degree \(D\).

Write \(W_m(q)=\sum_{t=0}^Dc_tq^t\), put \(c_{-1}=0\), and let
\(T=\lfloor D/2\rfloor\).

\begin{lemma}\label{lem:difference}
For \(0\leq t\leq T\),
\begin{align*}
c_t-c_{t-1}={}&p(t)-p(t-a)-p(t-b)+p(t-(2a+b))\\
&+p(t-(a+3b))-p(t-(2a+3b)). \tag{4}
\end{align*}
\end{lemma}

\begin{proof}
From (2),
\[
(1-q)W_m(q)=
\frac{(1-q^a)(1-q^b)(1-q^{a+b})(1-q^{a+2b})
(1-q^{2a+3b})}
{(1-q)(1-q^2)(1-q^3)(1-q^5)}.
\]
With \(X=q^a\) and \(Y=q^b\), its numerator is
\begin{align*}
&(1-X)(1-Y)(1-XY)(1-XY^2)(1-X^2Y^3)\\
={}&1-Y-X+X^2Y+XY^3-X^2Y^3+X^3Y^4-X^4Y^4\\
&\quad-X^3Y^6+X^4Y^7+X^5Y^6-X^5Y^7.
\end{align*}
The first six monomials give the shifts in (4).  Every remaining shift is at
least \(3a+4b\), while (3) gives
\(3a+4b>T\).  Hence the remaining terms do not contribute at or below the
midpoint.
\end{proof}

\section{Uniform positivity}

Let \(g(t)\) denote the right side of (4).  Assume first \(m\geq8\), so
\(a\geq34\), and write \(d=b-a\geq1\).  On each interval cut out by the five
positive shifts in (4), replace a positive translate \(p(x)\) by
\(Q(x)+91/360\), and a negative translate by \(Q(x)+1\).  Call the resulting
lower bound \(H(t)\).  By Lemma~\ref{lem:kernel}, \(g(t)\geq H(t)\).

Direct differentiation gives the following bounds.  Each entry in the last
column is \(360\) times a lower bound for \(H\) on the stated interval.
\[
\begin{array}{ccl}
\toprule
\text{interval}&\text{behavior of }H&360\min H\text{ is at least}\\
\midrule
{}[0,a)&\text{increasing}&91\\
{}[a,b)&\text{increasing}&2a^3+33a^2+162a-269\\
{}[b,2a+b)&\text{increasing}&
2a^3+6a^2d+33a^2+6ad^2+66ad+162a-629\\
{}[2a+b,a+3b)&\text{concave}&
22a^3+30a^2d-33a^2+6ad^2-66ad-162a-538\\
{}[a+3b,2a+3b)&\text{decreasing}&3(2ab(a+b-11)-149)\\
{}[2a+3b,T]&\text{decreasing}&3(2ab-269)\\
\bottomrule
\end{array} \tag{5}
\]
For completeness, on the third interval \(H'\) is a concave quadratic and
is positive at both endpoints.  On the fourth interval \(H\) is concave and
the smaller endpoint value is the one displayed.  On the fifth interval
\(H'\) is convex and negative at both endpoints.  On the last interval,
\(360H'=-12ab\); evaluation at the continuous midpoint
\((5a+7b-12)/2\) yields the final entry of (5), so the integer midpoint gives
at least that value.

Every expression in (5) is positive for \(a\geq34\) and \(d\geq1\).  For
example, the only non-obvious fourth entry can be grouped as
\[
(22a^3-33a^2-162a-538)+6ad^2+6ad(5a-11)>0.
\]
Thus all midpoint differences are positive for \(m\geq8\).

For \(1\leq m\leq7\), exact evaluation of (4) gives
\[
\begin{array}{c|rrrrrrr}
m&1&2&3&4&5&6&7\\ \hline
\min_{0\leq t\leq T}g(t)&0&0&0&1&1&1&1
\end{array}
\]
All midpoint differences are therefore nonnegative for every \(m\geq1\).
Symmetry proves Theorem~\ref{thm:main}.

\section{Reproducibility and claim boundary}

The proof uses exact arithmetic.  The companion checker verifies the
period-30 kernel, the inclusion--exclusion identity, all finite cases, and
the uniform lower-envelope calculations.  From the repository root run
\begin{verbatim}
python3 proof/qfib_width5_unimodality_proof.py
\end{verbatim}
The residue-class comparison in the introduction is checked exactly by
\begin{verbatim}
python3 experiments/multi_spacer_adversarial_and_width5_overlap.py
\end{verbatim}

The theorem proves only \(\qfib{m+5}{5}\) for \(m\geq1\).  It does not prove
the full two-parameter conjecture, width six, log-concavity, or a
combinatorial chain decomposition.

\paragraph{AI assistance.}
The author used AI assistance in developing the proof presentation and its
verification code.

\begin{thebibliography}{3}
\bibitem{BCK}
Nantel Bergeron, Cesar Ceballos, and Josef K\"ustner.
\newblock Elliptic and \(q\)-analogs of the Fibonomial numbers.
\newblock \emph{SIGMA}, 16:076, 2020.

\bibitem{CIMSY}
Brendan B. Connelly, Ezekiel Ito, Thomas C. Martinez, Olha Shevchenko, and
Kacey Yang.
\newblock Unimodality of \(q\)-Fibonomial coefficients for small cases.
\newblock \emph{arXiv:2605.12822v1 [math.CO]}, 2026.

\bibitem{ZhaoWidthFour}
Hainan Zhao.
\newblock The width-four \(q\)-Fibonomial coefficients are unimodal.
\newblock Zenodo, 2026. \href{https://doi.org/10.5281/zenodo.21826970}
{doi:10.5281/zenodo.21826970}.
\end{thebibliography}

\end{document}
